\documentclass[11pt,a4paper]{article}

\usepackage[utf8]{inputenc}
\usepackage[greek,english]{babel}
\usepackage[T1]{fontenc}
\usepackage{lmodern}
\usepackage{geometry}
\usepackage{booktabs}
\usepackage{hyperref}
\usepackage{graphicx}
\usepackage{amsmath}
\usepackage{listings}
\usepackage{caption}
\usepackage{subcaption}

\geometry{margin=2.5cm}
\hypersetup{colorlinks=true, linkcolor=blue, urlcolor=blue}

\title{%
  \textbf{Εξαμηνιαία Εργασία}\\[0.3em]
  Διαχείριση Δεδομένων Μεγάλης Κλίμακας\\[0.5em]
  \large DSML 2026 --- Ανάλυση εγκληματικότητας Los Angeles με Apache Spark
}
\author{Σταύρος Κοροβέσης \\ Α.Μ.: 03400297 \\ Username εργαστηρίου: dsml00297}
\date{Ακαδημαϊκό έτος 2025--26}

\begin{document}
\maketitle
\tableofcontents
\newpage

%==============================================================================
\section{Εισαγωγή}
%==============================================================================

Η παρούσα εργασία υλοποιεί τέσσερα analytical queries πάνω σε δημόσια δεδομένα
εγκληματικότητας και απογραφής του Los Angeles, με εκτέλεση στο Kubernetes cluster
του εργαστηρίου (cslab) μέσω \texttt{spark-submit}. Τα datasets βρίσκονται στο
HDFS path \texttt{/data/}. Ο κώδικας είναι διαθέσιμος στο αποθετήριο:
\texttt{<συμπληρώστε GitHub URL>}.

\textbf{Μέθοδος μέτρησης χρόνου:} όπως προτείνει ο πρακτικός οδηγός του μαθήματος,
ο χρόνος μετράται εντός του Python με \texttt{perf\_counter()} γύρω από το τελικό
Spark action (\texttt{write}). Η τιμή εξάγεται από τα driver logs ως
\texttt{QUERY\_ELAPSED\_SECONDS}. Για τις υλοποιήσεις RDD όπου το \texttt{collect()}
εκτελείται πριν τη μέτρηση, οι χρόνοι υποεκτιμούν το συνολικό κόστος·
σημειώνεται ρητά στην ενότητα~\ref{sec:q1}.

%==============================================================================
\section{Ζητούμενο 1: Σύγκριση μορφών αποθήκευσης (CSV vs Parquet)}
%==============================================================================

\subsection{Θεωρητική σύγκριση}

Το CSV είναι αναγνώσιμο και συμβατό, αλλά απαιτεί parsing, δεν ενσωματώνει schema
και είναι αναποτελεσματικό σε συμπίεση. Το \textbf{Parquet} είναι columnar format
με ενσωματωμένο schema, συμπίεση και predicate pushdown, κατάλληλο για analytics
στο Spark. Εναλλακτικά: JSON (ευέλικτο, μεγάλο), ORC (παρόμοιο με Parquet),
Avro (schema evolution).

\subsection{Μετατροπή}

Τα αρχεία \texttt{LA\_Crime\_Data\_2010\_2019.csv} και
\texttt{LA\_Crime\_Data\_2020\_2025.csv} μετατράπηκαν σε Parquet με το script
\texttt{convert\_parquet.py}. Χρόνος μετατροπής: \textbf{295.0\,s}.

\subsection{Benchmark στο Query 2}

\begin{table}[h]
\centering
\caption{Χρόνος εκτέλεσης Q2 (DataFrame) --- CSV vs Parquet}
\begin{tabular}{lcc}
\toprule
Format & Χρόνος (s) & Επιτάχυνση \\
\midrule
CSV    & 32.750 & 1.00$\times$ \\
Parquet & 4.247 & 7.72$\times$ \\
\bottomrule
\end{tabular}
\end{table}

Το Parquet είναι $\sim$8$\times$ ταχύτερο στο ίδιο query, λόγω μικρότερου I/O,
column pruning και αποφυγής CSV parsing. Για μεγάλα batch analytics workloads,
η μετάβαση σε Parquet (ή ORC) είναι συνήθως η βέλτιστη επιλογή.

%==============================================================================
\section{Ζητούμενο 2: Query 1 --- Μερίδιο εγκλημάτων σε δρόμο ανά τμήμα ημέρας}
\label{sec:q1}
%==============================================================================

\subsection{Περιγραφή}

Ταξινόμηση των τμημάτων ημέρας (Πρωί, Απόγευμα, Βράδυ, Νύχτα) κατά φθίνουσο
ποσοστό εγκλημάτων με \texttt{Premis Desc = "STREET"}.

\subsection{Config}
\texttt{2 executors $\times$ 1 core $\times$ 2GB}

\subsection{Αποτελέσματα}

\begin{table}[h]
\centering
\caption{Q1: Ποσοστό εγκλημάτων σε δρόμο ανά περίοδο}
\begin{tabular}{lrrr}
\toprule
Περίοδος & street\_pct (\%) & Σύνολο εγκλημάτων & Σε δρόμο \\
\midrule
Νύχτα     & 28.854 & 870235 & 251094 \\
Βράδυ     & 27.552 & 719695 & 198292 \\
Πρωί      & 18.881 & 693099 & 130866 \\
Απόγευμα  & 18.294 & 855099 & 156432 \\
\bottomrule
\end{tabular}
\end{table}

Η Νύχτα (21:00--04:59) και το Βράδυ (17:00--20:59) εμφανίζουν το υψηλότερο
ποσοστό εγκλημάτων σε δρόμο.

\subsection{Σύγκριση υλοποιήσεων}

\begin{table}[h]
\centering
\caption{Q1: Χρόνοι εκτέλεσης (1 run)}
\begin{tabular}{lcc}
\toprule
API & Script & Χρόνος (s) \\
\midrule
RDD & \texttt{q1\_rdd.py} & 1.92 \\
DataFrame & \texttt{q1\_df.py} & 21.62 \\
DataFrame + UDF & \texttt{q1\_df\_udf.py} & 224.56 \\
\bottomrule
\end{tabular}
\end{table}

Το DataFrame API είναι ταχύτερο από το UDF (224\,s), καθώς το UDF εμποδίζει
vectorized execution. Ο χρόνος RDD φαίνεται αποδεκτά χαμηλός επειδή η μέτρηση
ξεκινά μετά το \texttt{collect()}· το Catalyst optimizer στο DataFrame
εκτελεί aggregations πιο αποδοτικά στο cluster.

%==============================================================================
\section{Ζητούμενο 3: Query 2 --- Top-3 μήνες εγκλημάτων ανά έτος}
%==============================================================================

\subsection{Config}
\texttt{4 executors $\times$ 1 core $\times$ 2GB}

\subsection{Χρόνοι}

\begin{table}[h]
\centering
\caption{Q2: DataFrame vs Spark SQL}
\begin{tabular}{lcc}
\toprule
Υλοποίηση & Script & Χρόνος (s) \\
\midrule
DataFrame API & \texttt{q2\_df.py} & 327.66 \\
Spark SQL & \texttt{q2\_sql.py} & 101.97 \\
\bottomrule
\end{tabular}
\end{table}

Και οι δύο υλοποιήσεις εκφράζουν το ίδιο λογικό σχέδιο (group-by έτος/μήνας,
\texttt{RANK()} ανά έτος, φίλτρο top-3). Το Spark SQL είναι $\sim$3.2$\times$
ταχύτερο, πιθανότατα λόγω βελτιστοποίησης του Catalyst σε window aggregations
και λιγότερου Python overhead στο driver.

%==============================================================================
\section{Ζητούμενο 4: Query 3 --- Κατακεφαλήν εισόδημα ανά Zip Code}
%==============================================================================

\subsection{Ορισμός}

\[
\text{per\_capita} = \frac{\text{median\_household\_income}_{2021}}{\sum \text{POP20}_{\text{ανά ZCTA20}}}
\]

\subsection{Config}
\texttt{3 executors $\times$ 1 core $\times$ 2GB}

\subsection{Χρόνοι}

\begin{table}[h]
\centering
\caption{Q3: DataFrame vs RDD}
\begin{tabular}{lcc}
\toprule
Υλοποίηση & Script & Χρόνος (s) \\
\midrule
DataFrame & \texttt{q3\_df.py} & 285.30 \\
RDD & \texttt{q3\_rdd.py} & 3.39 \\
\bottomrule
\end{tabular}
\end{table}

Το DataFrame περιλαμβάνει ανάγνωση GeoJSON (explode features) και join με income.
Το RDD script μετρά μόνο το write μετά από in-memory \texttt{collect()}·
πραγματικός χρόνος RDD είναι σημαντικά μεγαλύτερος. Το DataFrame API είναι
προτιμητέο για joins και aggregations σε structured data.

%==============================================================================
\section{Ζητούμενο 5: Query 4 --- Πλησιέστερο αστυνομικό τμήμα}
%==============================================================================

\subsection{Μέθοδος}

Cross join εγκλημάτων (LAT/LON) με 21 αστυνομικά τμήματα (broadcast),
Haversine απόσταση, επιλογή ελάχιστου ανά έγκλημα, aggregation ανά \texttt{DIVISION}.

\subsection{Κάθετη κλιμάκωση (2 executors)}

\begin{table}[h]
\centering
\caption{Q4: Αύξηση cores/memory ανά executor}
\begin{tabular}{ccc}
\toprule
Cores/exec & Memory & Χρόνος (s) \\
\midrule
1 & 2GB & 295.14 \\
2 & 4GB & 356.36 \\
4 & 8GB & 89.38 \\
\bottomrule
\end{tabular}
\end{table}

Με 4 cores $\times$ 8GB ανά executor παρατηρείται σημαντική βελτίωση ($\sim$3.3$\times$
έναντι 1c/2g), καθώς το cross join και οι UDF distances ωφελούνται από
περισσότερους παράλληλους workers εντός κάθε executor.

\subsection{Οριζόντια κλιμάκωση (8 cores, 16GB συνολικά)}

\begin{table}[h]
\centering
\caption{Q4: Σταθερός συνολικός πόρος, διαφορετικός αριθμός executors}
\begin{tabular}{ccc}
\toprule
Config & Χρόνος (s) & Σχόλιο \\
\midrule
2 exec $\times$ (4c, 8GB) & 218.82 & λιγότερο scheduling overhead \\
4 exec $\times$ (2c, 4GB) & \textbf{176.56} & \textbf{καλύτερο} \\
8 exec $\times$ (1c, 2GB) & 230.81 & περισσότερο network shuffle \\
\bottomrule
\end{tabular}
\end{table}

Η διάταξη 4$\times$(2 cores, 4GB) δίνει το καλύτερο αποτέλεσμα: ισορροπία
μεταξύ παραλληλισμού και μεγέθους JVM heap ανά executor.

%==============================================================================
\section{Ζητούμενο 6: Join strategies (hints \& explain)}
%==============================================================================

\subsection{Query 3 (DataFrame join population--income)}

\begin{table}[h]
\centering
\caption{Q3: Επίδραση join hints}
\begin{tabular}{llc}
\toprule
Hint & Φυσικό plan & Χρόνος (s) \\
\midrule
broadcast & BroadcastHashJoin & 41.61 \\
merge & SortMergeJoin & \textbf{18.72} \\
shuffle\_hash & ShuffledHashJoin & 47.86 \\
shuffle\_replicate\_nl & BroadcastNestedLoopJoin & 81.77 \\
\bottomrule
\end{tabular}
\end{table}

Ο πίνακας income είναι μικρός ($\sim$300 zip codes). Το \texttt{merge}
(SortMergeJoin) αποδίδει καλύτερα όταν και οι δύο πλευρές είναι ήδη
μεγάλες μετά το aggregation των census blocks.

\subsection{Query 4 (cross join crimes $\times$ stations)}

\begin{table}[h]
\centering
\caption{Q4: Join hints σε cross join}
\begin{tabular}{llc}
\toprule
Hint & Φυσικό plan & Χρόνος (s) \\
\midrule
broadcast & BroadcastNestedLoopJoin & 294.64 \\
merge & BroadcastNestedLoopJoin & 294.91 \\
shuffle\_hash & BroadcastNestedLoopJoin & 282.84 \\
shuffle\_replicate\_nl & BroadcastNestedLoopJoin & \textbf{271.35} \\
\bottomrule
\end{tabular}
\end{table}

Στο cross join με 21 stations, ο optimizer επιλέγει \textbf{BroadcastNestedLoopJoin}
(broadcast του μικρού πίνακα stations). Η διαφορά μεταξύ hints είναι μικρή
($\sim$8\%) --- το bottleneck είναι ο υπολογισμός αποστάσεων σε $\sim$3M εγκλημάτων.

%==============================================================================
\section{Συμπεράσματα}
%==============================================================================

\begin{itemize}
  \item Parquet υπερέχει σαφώς του CSV για analytics (Q2: 7.7$\times$).
  \item DataFrame/SQL προτιμώνται έναντι RDD/UDF για complex queries.
  \item Στο Q4, η οριζόντια κλιμάκωση με 4 executors $\times$ 2 cores είναι βέλτιστη.
  \item Join hints έχουν ουσιαστικό αντίκτυπο στο Q3, όχι στο cross join του Q4.
\end{itemize}

%==============================================================================
\appendix
\section{Εντολές αναπαραγωγής}
%==============================================================================

\begin{lstlisting}[language=bash,basicstyle=\ttfamily\small]
source ~/bigdata-env.sh
cd ~/dsml2026-project/code

# Παράδειγμα Q1
spark-submit --py-files common.py \
  --conf spark.executor.instances=2 \
  --conf spark.executor.cores=1 \
  --conf spark.executor.memory=2g \
  q1_df.py \
  --output-base hdfs://hdfs-namenode.default.svc.cluster.local:9000/user/$DSML_USER/project-output

# Συλλογή χρόνων
bash ~/dsml2026-project/scripts/collect_timings.sh
\end{lstlisting}

\section{Δήλωση χρήσης LLM}

Βλέπε αρχείο \texttt{LLM\_USAGE.md} στο repository.

\end{document}
